\chapter{Evaluation}
\label{Evaluation}

This chapter reports the experimental results, structured by the three research questions defined in Chapter~\ref{Introduction}. Section~\ref{sec:eval-setup} describes the experimental setup. Section~\ref{sec:rq1} evaluates the impact of object re-identification (RQ1). Section~\ref{sec:rq2} compares the audio models across window and hop configurations (RQ2). Section~\ref{sec:rq3} evaluates the multimodal union condition against both unimodal baselines (RQ3). Section~\ref{sec:eval-summary} summarises the findings.

\section{Evaluation Setup}
\label{sec:eval-setup}

The evaluation uses a dataset of 30 synthetic surveillance scenes (10 s, 1280$\times$720, 24 fps) generated with Dreamina Seedance. Each scene is annotated with per-frame ground truth for visual relations and sound classes, and event-level binary verdicts. The 30 scenes span 6 event types with 5 positive instances each.

Four VLM models were evaluated on condition A: Pixtral 12B, Ministral 3-14B, Gemini 2.5 Flash, and Gemini 3.6 Flash. Two audio providers were evaluated on condition B: PANNs CNN14 and Qwen2-Audio-7B-Instruct, each across 8 window/hop configurations. Condition C combined every VLM with every audio provider via union. A re-identification ablation was conducted for each VLM by running without the tracking prompt.

The primary metric is the F1 score (\S\ref{sec:metrics}), reported alongside precision, recall, and per-event TP/FP/FN counts.

\section{RQ1: Object Re-Identification}
\label{sec:rq1}

The first research question asks whether object re-identification is necessary for visual event detection. Table~\ref{tab:reid_ablation} reports the F1, precision, recall, and per-event counts for each VLM on all 30 scenes, with and without re-ID, and the per-metric deltas ($\Delta$) when re-ID is applied.

\begin{table}[h]
  \centering
  \caption{Re-identification ablation: all four VLM models on 30 scenes, with and without re-ID. Deltas ($\Delta$) show the change when re-ID is applied.}
  \label{tab:reid_ablation}
  \resizebox{\textwidth}{!}{%
  \begin{tabular}{l ccc cc cc ccc cc cc}
    \toprule
    VLM & \multicolumn{3}{c}{F1} & \multicolumn{2}{c}{Precision} & \multicolumn{2}{c}{Recall} & \multicolumn{3}{c}{TP} & \multicolumn{2}{c}{FP} & \multicolumn{2}{c}{FN} \\
    \cmidrule(lr){2-4} \cmidrule(lr){5-6} \cmidrule(lr){7-8} \cmidrule(lr){9-11} \cmidrule(lr){12-13} \cmidrule(lr){14-15}
         & no-ID & re-ID & $\Delta$ & no-ID & re-ID & no-ID & re-ID & no-ID & re-ID & $\Delta$ & no-ID & re-ID & no-ID & re-ID \\
    \midrule
    \textbf{Gemini 3.6 Flash} & \textbf{0.720} & \textbf{0.842} & \textbf{+0.122} & \textbf{0.900} & \textbf{0.889} & \textbf{0.600} & \textbf{0.800} & \textbf{18} & \textbf{24} & \textbf{+6} & \textbf{2} & \textbf{3} & \textbf{12} & \textbf{6} \\
    \midrule
    Ministral 3-14B & 0.691 & 0.750 & +0.059 & 0.760 & 0.808 & 0.633 & 0.700 & 19 & 21 & +2 & 6 & 5 & 11 & 9 \\
    \midrule
    Gemini 2.5 Flash & 0.571 & 0.653 & +0.082 & 0.737 & 0.842 & 0.467 & 0.533 & 14 & 16 & +2 & 5 & 3 & 16 & 14 \\
    \midrule
    Pixtral 12B & 0.712 & 0.679 & $-$0.033 & 0.724 & 0.783 & 0.700 & 0.600 & 21 & 18 & $-$3 & 8 & 5 & 9 & 12 \\
    \bottomrule
  \end{tabular}%
  }
\end{table}

Re-identification consistently improves recall by recovering the tracking-dependent events. For Gemini 3.6 Flash, recall rises from 0.600 to 0.800, driven by loitering (0/5 to 4/5), vehicle escape (0/5 to 1/5), and handoff (4/5 to 5/5). Precision also increases for three of the four models because single-frame fragments that produce spurious detections without re-ID are suppressed by the tracking prompt: Gemini 2.5 Flash shows the largest precision gain (0.737 to 0.842). False positives decrease for Ministral 3-14B (6 to 5) and Pixtral 12B (8 to 5), while increasing slightly for Gemini 3.6 Flash (2 to 3) primarily in the fight event where the VLM detects physical altercation in non-fight scenes, an imperfection unrelated to re-ID. The net F1 improvement is substantial for three models (Gemini 3.6 Flash +0.122, Gemini 2.5 Flash +0.082, Ministral 3-14B +0.059). Pixtral 12B is the exception: it detects more events without re-ID (21 vs.\ 18 TP), though with markedly lower precision (0.724 vs.\ 0.783), suggesting its raw predictions fire duration-gated queries even without consistent identity tracking. The effect is concentrated on events requiring persistent object identity; short-duration events such as fight and vehicle collision are largely unaffected in all four models.

\section{RQ2: Audio Model Ablation}
\label{sec:rq2}

The second research question compares Qwen2-Audio-7B-Instruct (LALM) against PANNs CNN14 across sliding window and hop configurations. Table~\ref{tab:audio_ablation} shows the full ablation on the 20 audio-relevant scenes (fight, gunshot, vehicle escape, and vehicle collision have acoustic components; handoff and loitering do not).

\begin{table}[h]
  \centering
  \caption{Window/hop ablation for audio-only (condition B) across 20 scenes.}
  \label{tab:audio_ablation}
  \resizebox{\textwidth}{!}{%
  \begin{tabular}{lcccccccc}
    Model & Window (s) & Hop (s) & F1 & Precision & Recall & TP & FP & FN \\
    \midrule
    \textbf{Qwen2-Audio} & \textbf{5.0} & \textbf{2.50} & \textbf{0.788} & \textbf{1.000} & \textbf{0.650} & \textbf{13} & \textbf{0} & \textbf{7} \\
    Qwen2-Audio & 2.5 & 1.25 & 0.750 & 1.000 & 0.600 & 12 & 0 & 8 \\
    Qwen2-Audio & 5.0 & 5.00 & 0.667 & 1.000 & 0.500 & 10 & 0 & 10 \\
    Qwen2-Audio & 2.5 & 2.50 & 0.621 & 1.000 & 0.450 & 9 & 0 & 11 \\
    Qwen2-Audio & 1.0 & 1.00 & 0.389 & 0.438 & 0.350 & 7 & 9 & 13 \\
    Qwen2-Audio & 1.0 & 0.50 & 0.368 & 0.389 & 0.350 & 7 & 11 & 13 \\
    Qwen2-Audio & 10.0 & 10.00 & 0.320 & 0.800 & 0.200 & 4 & 1 & 16 \\
    Qwen2-Audio & 10.0 & 5.00 & 0.320 & 0.800 & 0.200 & 4 & 1 & 16 \\
    \midrule
    PANNs CNN14 & 2.5 & 1.25 & 0.429 & 0.750 & 0.300 & 6 & 2 & 14 \\
    PANNs CNN14 & 5.0 & 5.00 & 0.400 & 1.000 & 0.250 & 5 & 0 & 15 \\
    PANNs CNN14 & 1.0 & 0.50 & 0.385 & 0.833 & 0.250 & 5 & 1 & 15 \\
    PANNs CNN14 & 2.5 & 2.50 & 0.370 & 0.714 & 0.250 & 5 & 2 & 15 \\
    PANNs CNN14 & 5.0 & 2.50 & 0.370 & 0.714 & 0.250 & 5 & 2 & 15 \\
    PANNs CNN14 & 1.0 & 1.00 & 0.320 & 0.800 & 0.200 & 4 & 1 & 16 \\
    PANNs CNN14 & 10.0 & 10.00 & 0.261 & 1.000 & 0.150 & 3 & 0 & 17 \\
    PANNs CNN14 & 10.0 & 5.00 & 0.261 & 1.000 & 0.150 & 3 & 0 & 17 \\
    \bottomrule
  \end{tabular}%
  }
\end{table}

Several patterns emerge from the ablation. First, Qwen2-Audio achieves \textbf{perfect precision} (1.000) for all configurations with a window of 2.5 s or larger: every detection is correct. Its best micro F1 (0.788) occurs at (5.0 s, 2.5 s), with the 2.5 s / 1.25 s configuration close behind (0.750). Smaller windows (1.0 s) degrade precision because the model has insufficient acoustic context; larger windows (10.0 s) miss short-duration events. Second, PANNs CNN14 achieves its best F1 of 0.429 at (2.5 s, 1.25 s) but never exceeds this value. PANNs performs well on gunshot detection (high recall) but suffers from false positives and low recall on other events. The fight audio query uses a BEFORE($\delta=120$), matching the ISEQL paper definition, though it did not improve recall because Qwen2-Audio detects impact before shout in the two relevant scenes. The 2.5 s / 1.25 s configuration balances context size and temporal resolution and is used in the multimodal experiments; at this configuration Qwen2-Audio still achieves perfect precision with F1 = 0.750.

\section{RQ3: Multimodal vs.\ Unimodal}
\label{sec:rq3}

\subsection{Visual-Only Results}

Table~\ref{tab:vlm_comparison} reports the visual-only (condition A) results. All models use object re-identification.

\begin{table}[h]
  \centering
  \caption{Visual-only (condition A) across four VLM models on 30 scenes.}
  \label{tab:vlm_comparison}
  \begin{tabular}{lcccccc}
    \toprule
    VLM & F1 & Precision & Recall & TP & FP & FN \\
    \midrule
    \textbf{Gemini 3.6 Flash} & \textbf{0.842} & \textbf{0.889} & \textbf{0.800} & \textbf{24} & \textbf{3} & \textbf{6} \\
    Ministral 3-14B & 0.750 & 0.808 & 0.700 & 21 & 5 & 9 \\
    Pixtral 12B & 0.679 & 0.783 & 0.600 & 18 & 5 & 12 \\
    Gemini 2.5 Flash & 0.653 & 0.842 & 0.533 & 16 & 3 & 14 \\
    \bottomrule
  \end{tabular}
\end{table}

Gemini 3.6 Flash achieves the highest F1 (0.842), detecting 24 of 30 events. The gap between the best and worst model is 0.189 F1 points, with the largest differences on events requiring accurate re-identification across multiple frames (vehicle escape, handoff). The improvement over previous results is due to fixing the explosion detection query: explosion scenes involve vehicles and objects, not persons, so the earlier query silently dropped all explosion detections.

\subsection{Multimodal Union Results}

Table~\ref{tab:multimodal_results} shows the multimodal (condition C) results combining each VLM with each audio provider via union. Each pair is reported at its best window/hop configuration; when several configurations achieve the same F1, all tied configurations are shown.

\begin{table}[h]
  \centering
  \caption{Multimodal (condition C) across all VLM and audio combinations, best window/hop per pair.}
  \label{tab:multimodal_results}
  \resizebox{\textwidth}{!}{%
  \begin{tabular}{lccccccc}
    VLM + Audio Model & Window/Hop (s) & F1 & Precision & Recall & TP & FP & FN \\
    \midrule
    \textbf{Gemini 3.6 Flash + Qwen2-Audio} & \textbf{2.5/1.25, 5.0/2.5, 5.0/5.0} & \textbf{0.935} & \textbf{0.906} & \textbf{0.967} & \textbf{29} & \textbf{3} & \textbf{1} \\
    Ministral 3-14B + Qwen2-Audio & 5.0/2.5, 5.0/5.0 & 0.889 & 0.848 & 0.933 & 28 & 5 & 2 \\
    Gemini 3.6 Flash + PANNs & 5.0/5.0, 10/5, 10/10 & 0.862 & 0.893 & 0.833 & 25 & 3 & 5 \\
    Pixtral 12B + Qwen2-Audio & 5.0/2.5, 5.0/5.0 & 0.833 & 0.833 & 0.833 & 25 & 5 & 5 \\
    Gemini 2.5 Flash + Qwen2-Audio & 5.0/2.5 & 0.821 & 0.885 & 0.767 & 23 & 3 & 7 \\
    Ministral 3-14B + PANNs & 5.0/5.0, 10/5, 10/10 & 0.772 & 0.815 & 0.733 & 22 & 5 & 8 \\
    Gemini 2.5 Flash + PANNs & 5.0/5.0 & 0.706 & 0.857 & 0.600 & 18 & 3 & 12 \\
    Pixtral 12B + PANNs & 5.0/5.0, 10/5, 10/10 & 0.704 & 0.792 & 0.633 & 19 & 5 & 11 \\
    \bottomrule
  \end{tabular}%
  }
\end{table}

The best combination achieves F1 = 0.935, detecting 29 of 30 events with 3 false positives and 1 false negative. This result is robust across the window/hop configuration: Gemini 3.6 Flash + Qwen2-Audio reaches F1 = 0.935 at the audio-optimal 5.0 s / 2.5 s configuration as well as at 2.5 s / 1.25 s and 5.0 s / 5.0 s. The top combinations all use Qwen2-Audio, demonstrating that the LALM's perfect precision carries through to the multimodal setting: false positives from audio would propagate to condition C, so precise audio models are essential. The complete ranked ablation across all 64 combinations (4 VLM models $\times$ 2 audio models $\times$ 8 window/hop configurations) is reported in Appendix~\ref{AppendixA} (Table~\ref{tab:full_ablation}).

\subsection{Thesis Claim Verification}

The thesis claim is $C \ge \max(A, B)$ ,  the multimodal condition never performs worse than the better of the two unimodal conditions. Table~\ref{tab:thesis_claim} verifies this claim per event for the best combination.

\begin{table}[h]
  \centering
  \caption{Per-event breakdown: Gemini 3.6 Flash + Qwen2-Audio.}
  \label{tab:thesis_claim}
  \begin{tabular}{lcccc}
    \toprule
    Event & Scenes & A (visual) & B (audio) & C (union) \\
    \midrule
    Fight & 5 & 5/5 & 0/5 & 5/5 \\
    Gunshot/Explosion & 5 & 4/5 & 5/5 & 5/5 \\
    Vehicle Escape & 5 & 1/5 & 4/5 & 5/5 \\
    Vehicle Collision & 5 & 5/5 & 3/5 & 5/5 \\
    Loitering & 5 & 4/5 & -- & 4/5 \\
    Handoff & 5 & 5/5 & -- & 5/5 \\
    \midrule
    Overall & 30 & 24/30 & 12/20 & 29/30 \\
    \bottomrule
  \end{tabular}
\end{table}

The claim holds on all 30 scenes: condition C is never worse than $\max(A, B)$. The union is most beneficial when the modalities are complementary: audio contributes 4 additional true positives on vehicle escape that visual alone would have missed (visual 1/5, audio 4/5, union 5/5), while visual carries the fight detection (5/5) that audio completely misses. The LALM detects all 5 gunshot/explosion events through synonym resolution, adding one detection over visual alone. The two visual-only events (handoff, loitering) are unchanged in C since they have no audio query; handoff is recovered to 5/5 by re-identification.

\section{Summary}
\label{sec:eval-summary}

The evaluation establishes three findings. \textbf{RQ1:} Object re-identification is essential for tracking-dependent events; without it, loitering, vehicle escape, and handoff are largely missed (e.g. Gemini 3.6 Flash loitering 0/5 to 4/5, escape 0/5 to 1/5). \textbf{RQ2:} Qwen2-Audio achieves perfect precision (1.000) and a micro F1 of 0.788 at 5.0 s / 2.5 s (0.750 at 2.5 s / 1.25 s), outperforming PANNs (best F1 = 0.429) across all configurations. \textbf{RQ3:} The multimodal union achieves F1 = 0.935 with the best combination (Gemini 3.6 Flash + Qwen2-Audio, 29/30 TP), outperforming both unimodal baselines (visual F1 = 0.842, audio F1 = 0.788), and the claim $C \ge \max(A, B)$ holds on all 30 scenes.
